\section{Discussion}
\label{sec:discussion}
\para{What does this study establish?}
The strongest supported finding is that prompt context is a major confound in this instructed-probe evaluation. A pre-answer feature separates the observed false-output conditions perfectly, and known-correct pressure responses frequently cross a threshold calibrated in neutral contexts. Consequently, a high instructed-condition \auroc cannot by itself establish mechanism specificity. We do not infer that the response probe contains no belief signal: one recovered test example cannot distinguish that possibility from prompt and retrieval features. Our results complement concerns about transfer, belief verification, and calibration in prior work~\cite{goldowskydill2025detecting,kretschmar2025liars,cooney2026did,hopkins2026fine}; they do not refute metrics obtained on different models, labels, or settings.

\para{Why is mechanism prevalence unresolved?}
The behavioral partition deliberately leaves most errors unassigned. Low recovered yield could reflect weak fictional incentives, limited model capability, instruction hierarchy, anchoring during follow-up, or incomplete knowledge measurement. Neutral accuracy is indirect evidence about available beliefs; a private correction is an intervention that can change retrieval. The proxy also has a validity problem at the individual-example level: one recovered answer, ``The seeme,'' is an uninterpretable band-name fragment. Excluding it leaves only a Simferopol response followed by a Yalta correction. Neither example warrants a psychological-intent claim. The small operational share does not imply that deployed models rarely lie.

\para{Limits of model, data, and measurement.}
We evaluate one historical 7B model, one main generation per arm, and fictional stakes. Fixed arm order and a single full-study seed limit sampling conclusions. \triviaqa and \maskbench may have appeared in pretraining; canonical entity grouping cannot eliminate every alias or topic dependency. Reference errors, strict normalization, numeric tolerance, and fallible same-family model auditing affect labels. We do not perform independent multi-human adjudication or comprehensive source verification. Conditioning on falsehood is post-treatment selection, and the equal arm mixture is artificial. In particular, our uncertain-error category covers only neutral low-evidence falsehoods; it is not an exhaustive category for pressured epistemic errors.

\para{Limits of detector comparisons.}
Only one recovered test positive, zero recovered reputation positives, and 45 calibration responses preclude reliable verified-mechanism ranking or precise low-FPR claims. Small-cell bootstrap intervals cannot create evidence for unseen types. The probes are model-specific adaptations, not the original published checkpoints. Three normalized-string samples do not reproduce entailment-based semantic entropy, and our mean response features differ from original \sep token choices~\cite{kossen2024semantic}. We do not implement a trained follow-up detector, train deceptive organisms, or perform a cross-model sweep. External \maskbench results have severe abstention and a modified prompt; they supply a bounded calibration check only.

\para{Implications for monitoring and broader impact.}
Pressure-matched correct answers and prompt-only features should accompany claims about deception probes. In this study, they reveal a risk of penalizing answers for the context in which they were produced. Detector scores should therefore not serve as ground-truth labels for dishonesty or as stand-alone evidence for sanctions. At the same time, these results provide no assurance that a low-scoring answer is honest. Misuse could take the form of overconfident monitoring decisions or treating diagnostic weaknesses as proof that monitoring is futile. Our contribution is a reproducible measurement protocol and its observed limits, not an operational certification of model honesty.
